Every step below is given in full over $\C$. The cited real results are the same statements for real quadratic spaces; their arguments use only a field in which $2$ is invertible, and they are cited as the real case of each step, not as a source for the complex statement.

(i) Existence and uniqueness. Let $I_Q \subset T(V)$ be the two-sided ideal generated by the $v \otimes v - Q(v)$. A $\C$-linear $f : V \to A$ with $f(v)^{2} = Q(v)$ extends to a unique homomorphism $T(V) \to A$, which kills $I_Q$ and so factors uniquely through $T(V)/I_Q$; two universal pairs are isomorphic by the usual two-sided argument (real case: \claimref{std-algebra.existence-uniqueness}). Polarising $v^{2} = Q(v)$ gives $uv + vu = 2B(u, v)$ (\claimref{std-algebra.polarisation}, real case). Blade basis: words in a basis $e_1, \dots, e_n$ span, and $e_j e_i = -e_i e_j + 2B(e_i, e_j)$, $e_i e_i = Q(e_i)$ reduce every word to increasing ones by induction on length and inversions. For independence, let $c(v) = v \wedge{} + \iota_{B(v, \cdot)}$ on $\Lambda V$ (over $\C$); since $v \wedge v = 0$, $\iota^{2} = 0$ and $\iota_\alpha(v \wedge \omega) + v \wedge \iota_\alpha \omega = \alpha(v)\omega$, $c(v)^{2} = Q(v)$, so the universal property gives a representation, and $e_I \cdot 1 = e_{i_1} \wedge \dots \wedge e_{i_k} + (\text{lower degree})$ is triangular against the monomial basis of $\Lambda V$, so the $e_I$ are independent (real case: \claimref{std-algebra.blade-basis}). Nothing here uses nondegeneracy. The grade involution is the extension of $v \mapsto -v$ and reversion the extension of $v \mapsto v$ into the opposite algebra, both by the universal property; their signs $(-1)^{k}$ and $(-1)^{k(k-1)/2}$ on blades of grade $k$ count transpositions (real case: \claimref{std-algebra.involution-properties}). For a real $(V_\R, Q_\R)$, the $\R$-linear map $V_\R \to \Cl(V_\R) \otimes_\R \C$ extends $\C$-linearly to $V_\R \otimes \C$ with $(v \otimes z)^{2} = z^{2}Q(v)$, polarisation gives the relation for $Q_\C$ on all of $V_\R \otimes \C$, and the complex universal property gives a homomorphism $\Cl(V_\R \otimes \C) \to \Cl(V_\R) \otimes \C$, onto since the image contains $V_\R$ and $\mathrm{i}$, and bijective by the dimension $2^{n}$.

(ii) If $Q \ne 0$ there is a vector with $Q(e) \ne 0$: otherwise $B(u, v) = \tfrac12(Q(u + v) - Q(u) - Q(v)) = 0$, contradicting nondegeneracy when $n \ge 1$. Then $V = \C e \oplus e^{\perp}$ with $e^{\perp}$ nondegenerate of dimension $n - 1$, and induction gives an orthogonal basis with $Q(e_i) \ne 0$. Every nonzero complex number has a square root, so $e_i / \sqrt{Q(e_i)}$ is orthonormal. The isometry with $\C^{n}$ extends to an algebra isomorphism by the universal property.

(iii) A totally null $U$ lies in $U^{\perp}$, which has dimension $n - \dim U$ by nondegeneracy, so $\dim U \le k$. For $w_j = \tfrac12(e_{2j-1} + \mathrm{i} e_{2j})$, $Q(w_j) = \tfrac14(1 + \mathrm{i}^{2}) = 0$ and distinct $w_j$ involve disjoint orthogonal pairs, so $W$ is totally null of dimension $k$; likewise $v_j = \tfrac12(e_{2j-1} - \mathrm{i} e_{2j})$, with $B(v_i, w_j) = \tfrac14(1 - \mathrm{i}^{2})\delta_{ij} = \tfrac12\delta_{ij}$. For an arbitrary totally null $W$ with basis $w_1, \dots, w_k$: the map $x \mapsto B(x, \cdot)|_W$ from $V$ to $W^{*}$ has kernel $W^{\perp}$ of dimension $n - k$, so it is onto, and there are $x_j$ with $B(x_j, w_i) = \tfrac12\delta_{ij}$. Put $v_j = x_j - \sum_l B(x_j, x_l)\, w_l$. Then $B(v_j, w_i) = \tfrac12\delta_{ij}$ still, and $B(v_i, v_j) = B(x_i, x_j) - \tfrac12 B(x_j, x_i) - \tfrac12 B(x_i, x_j) = 0$. So $W' = \operatorname{span}\{v_j\}$ is totally null, $W \cap W' = 0$ by the nondegenerate pairing, and $V = W \oplus W'$ by dimension. By $uv + vu = 2B(u, v)$, $v_i w_j + w_j v_i = \delta_{ij}$ and the $v$'s and the $w$'s anticommute among themselves: the CAR (real case: \claimref{std-algebra.witt-basis-car}).
Fock model (real case: \claimref{std-algebra.fock-representation}). For $x = w + v$ the operator $c(x) = w \wedge{} + \iota_{2B(v, \cdot)}$ on $\Lambda W$ satisfies $c(x)^{2} = \{w \wedge, \iota_{2B(v, \cdot)}\} = 2B(v, w) = Q(x)$, because $W, W'$ are totally null; the universal property gives $\phi : \Cl(V, Q) \to \End_\C(\Lambda W)$. With $e_S = w_{s_1} \wedge \dots \wedge w_{s_m}$ and $\iota_{2B(v_i, \cdot)} w_j = \delta_{ij}$, $\phi(v_j w_j)$ fixes $e_S$ for $j \notin S$ and kills it for $j \in S$, so $\phi(f)$ is the projection onto $\C e_\emptyset$ along the other $e_S$. Then $\phi(w_I)\phi(f)\phi(v_{j_m} \cdots v_{j_1})$ is $\pm$ the elementary operator $e_J \mapsto e_I$, so $\phi$ is onto; both sides have dimension $4^{k}$ by (i), so $\phi$ is an isomorphism. As $\phi(f)$ has rank one, $f$ is a primitive idempotent, $\phi(v_i)\phi(f) = 0$, and $\Cl f \cong \End(\Lambda W)\phi(f) \cong \Lambda W$ by $w_I f \mapsto e_I$. Each vector shifts exterior degree by $\pm 1$, so even elements preserve $\Lambda^{\mathrm{even}}$ and $\Lambda^{\mathrm{odd}}$; $\phi$ maps $\Cl^{+}$ (dimension $2^{n-1} = 2 \cdot 4^{k-1}$) injectively into $\End(\Lambda^{\mathrm{even}}) \oplus \End(\Lambda^{\mathrm{odd}})$, of the same dimension, so onto.

(iv) In an orthonormal basis, moving each $e_i$ of the second factor of $(e_1 \cdots e_n)(e_1 \cdots e_n)$ to its partner takes $n(n-1)/2$ transpositions in all, so $(e_1 \cdots e_n)^{2} = (-1)^{n(n-1)/2}$ (real case: \claimref{primer.sign}), and with $m = n(n-1)/2$, $\omega_\C^{2} = \mathrm{i}^{2m}(-1)^{m} = 1$. Moving $e_j$ through $e_1 \cdots e_n$ passes $n - 1$ anticommuting factors and one commuting one, so $e_j \omega_\C = (-1)^{n-1}\omega_\C e_j = \omega_\C e_j$ for $n$ odd: $\omega_\C$ is central (real case: \claimref{primer.centre}). So $P_\pm$ are central with $P_\pm^{2} = P_\pm$, $P_+P_- = 0$, $P_+ + P_- = 1$, and $\Cl = \Cl P_+ \oplus \Cl P_-$ as algebras. Multiplication by the odd invertible $\omega_\C$ maps $\Cl^{+}$ injectively into $\Cl^{-}$, of the same dimension $2^{n-1}$, so $\Cl^{-} = \Cl^{+}\omega_\C$ and $\Cl P_\pm = \Cl^{+} P_\pm$ (as $\omega_\C P_\pm = \pm P_\pm$). The map $a \mapsto a P_\pm$ is a homomorphism on $\Cl^{+}$ ($P_\pm$ central idempotent), injective since $aP_\pm = \tfrac12(a \pm a\omega_\C)$ has even part $\tfrac12 a$. For $j < n$ put $g_j = \mathrm{i}\, e_j e_n \in \Cl^{+}$: $g_j^{2} = -e_j e_n e_j e_n = e_j^{2} e_n^{2} = 1$ and $g_j g_l = -g_l g_j$ for $j \ne l$, so the universal property gives $\Cl_{n-1}(\C) \to \Cl^{+}$; its image contains $g_I = \pm \mathrm{i}^{\lvert I\rvert} e_I e_n^{\lvert I \rvert \bmod 2}$, hence every even blade, so it is onto and, both sides of dimension $2^{n-1}$, an isomorphism (real case: \claimref{std-algebra.even-subalgebra-isomorphism}). With (iii), $\Cl_{2k+1}(\C) \cong \Mat_{2^{k}}(\C)^{\oplus 2}$, and $\Mat_{2^{k+1}}(\C) = \Mat_2(\Mat_{2^{k}}(\C))$ gives $\Cl_{n+2}(\C) \cong \Mat_2(\Cl_n(\C))$ for both parities. By (i) and (ii) these are the types of $\Cl(V) \otimes \C$ for real nondegenerate $V$ stated in part (i) of \claimref{std-algebra.complexification}, now with the complex algebra defined directly.

(v) (Real case: \claimref{std-algebra.chevalley-pairing}.) For monomials $x = e_S$, $y = e_T$, $\tau(e_S) \wedge e_T$ is $\pm e_{S \cup T}$ if $S \cap T = \emptyset$ and $0$ otherwise, so $\beta(e_S, e_T) = \pm 1$ exactly when $T$ is the complement of $S$: the Gram matrix is a signed permutation matrix, so $\beta$ is nondegenerate and pairs degree $m$ with $k - m$. Since $\tau$ reverses wedge products, $\tau(\tau x \wedge y) = \tau y \wedge x$, and $\tau$ acts on the top degree by $(-1)^{k(k-1)/2}$, so $\beta(y, x) = (-1)^{k(k-1)/2}\beta(x, y)$. For $w \in W$: $\tau(w \wedge x) = \tau x \wedge w$, so $\beta(wx, y) = \operatorname{top}(\tau x \wedge w \wedge y) = \beta(x, wy)$. For $v \in W'$ acting by the degree $-1$ derivation $\iota = \iota_{2B(v, \cdot)}$, and $x \in \Lambda^{m}$: $\iota(\tau x) = (-1)^{m+1}\tau(\iota x)$, and $\operatorname{top}\iota(\tau x \wedge y) = 0$ with $\iota(\tau x \wedge y) = \iota(\tau x) \wedge y + (-1)^{m}\tau x \wedge \iota y$, so $\beta(vx, y) = \beta(x, vy)$. Every vector of $V = W \oplus W'$ is therefore self-adjoint, and for $a = u_1 \cdots u_m$, $\beta(ax, y) = \beta(x, u_m \cdots u_1 y) = \beta(x, \tilde a y)$, extended linearly. All of this is $\C$-bilinear. A bivector has $\tilde A = -A$, so $\beta(Ax, y) = -\beta(x, Ay)$, and if $\tilde R R = 1$ then $\beta(Rx, Ry) = \beta(x, \tilde R R y) = \beta(x, y)$. For $k = 3$, $(-1)^{3} = -1$ gives antisymmetry, and $\beta$ pairs degree $m$ with $3 - m$, of opposite parity, so it vanishes on $\Lambda^{\mathrm{even}} \times \Lambda^{\mathrm{even}}$ and $\Lambda^{\mathrm{odd}} \times \Lambda^{\mathrm{odd}}$, and nondegeneracy makes the restriction to $\Lambda^{\mathrm{even}} \times \Lambda^{\mathrm{odd}}$ perfect. Finally $\Cl(2,4) \otimes \C \cong \Cl_6(\C)$ by (i) and (ii); the isotropic pair $w_1 = a_1 + b_1$, $w_2 = a_2 + b_2$, $w_3 = b_3 + \mathrm{i} b_4$ (with $Q(a_i) = 1$, $Q(b_i) = -1$) used in the conformal-spacetime results is such a $W$, completed by $v_j = \tfrac14(a_j - b_j)$ ($j = 1, 2$) and $v_3 = -\tfrac14(b_3 - \mathrm{i} b_4)$.
